
\documentclass[UTF8,a4paper,landscape]{ctexart}
\usepackage[margin=0.36cm,columnsep=0.24cm]{geometry}
\usepackage{amsmath,amssymb,mathtools}
\usepackage[version=4]{mhchem}
\usepackage{xcolor,multicol,enumitem}
\pagestyle{empty}
\setlength{\parindent}{0pt}
\setlength{\parskip}{0.2pt}
\setlength{\fboxsep}{1.1pt}
\setlength{\fboxrule}{0.32pt}
\setlength{\columnseprule}{0.18pt}
\setlength{\abovedisplayskip}{1pt}
\setlength{\belowdisplayskip}{1pt}
\setlength{\jot}{0.5pt}
\setlist[itemize]{nosep,leftmargin=0.9em,label=\raisebox{0.2ex}{\tiny$\bullet$}}
\setlist[enumerate]{nosep,leftmargin=1.1em}
\definecolor{head}{gray}{0}
\definecolor{sub}{gray}{0}
\definecolor{warn}{gray}{0}
\newcommand{\Mn}{\overline{M}_n}\newcommand{\Mw}{\overline{M}_w}\newcommand{\Mz}{\overline{M}_z}
\newcommand{\Xn}{\overline{X}_n}\newcommand{\Xw}{\overline{X}_w}\newcommand{\PDI}{\mathrm{PDI}}
\newcommand{\Ri}{R_i}\newcommand{\Rp}{R_p}\newcommand{\Rt}{R_t}\newcommand{\kp}{k_p}\newcommand{\kt}{k_t}\newcommand{\kd}{k_d}\newcommand{\dd}{\mathrm d}
\newcommand{\sect}[1]{\vspace{1.15pt}\noindent\fbox{\sffamily\bfseries #1}\par\vspace{0.35pt}}
\newcommand{\subsect}[1]{\vspace{0.5pt}{\sffamily\bfseries #1}\quad}
\newcommand{\key}[1]{{\sffamily\bfseries\underline{#1}}}
\newenvironment{dense}{\fontsize{8.35}{8.75}\selectfont\raggedright\sloppy\emergencystretch=8em}{ }
\begin{document}
\begin{center}
{\bfseries\large 合成化学II 两页压缩笔记}\quad {\small 覆盖作业与12页版核心内容}
\end{center}
\vspace{-2.0mm}
\begin{dense}
\begin{multicols}{4}
\sect{0. 符号、单位、总判断}
$[M],[I],[S],[P]$：单体、引发剂、溶剂/转移剂、聚合物浓度；$\Ri,\Rp,\Rt$：引发/聚合/终止速率；$f$：引发剂效率；$[M\cdot]$：增长自由基总浓度；$\nu$：动力学链长；$p,r$：反应程度、官能团当量比；共聚中 $f_i,F_i$ 为投料/瞬时共聚物组成，$r_i$ 为竞聚率；$\bar f,f_w,p_c$：平均/重均官能度、凝胶点。
\key{公式符号补表：}$i$为组分编号；$N_i,M_i$为第$i$类链的分子数、分子量，$M_0$为重复单元摩尔质量，$X,\Xn,\Xw$为单链/数均/重均聚合度，$N_0,N,N_x,x$为初始分子数、当前分子数、$x$聚体数、转化率；$\kd,k_i,\kp,\kt,k_{tr}$为分解/引发/增长/终止/转移速率常数，$C_X,C_M,C_I,C_S,C_P=k_{tr,X}/k_p$；$C,C_0,t,n,[Z],t_{ind}$为官能团浓度、初浓度、时间、阻聚计量数、阻聚剂浓度、诱导期；$[P^-],[I]_{eff}$为阴离子活性链和有效引发剂浓度；$\Delta G,\Delta H,\Delta S,T_c$为自由能、焓、熵、上限温度；$f_A,f_{B,w}$为A类官能度和B类重均官能度。$R$为通用自由基/引发剂残基，$P_n,P_m$为聚合度$n,m$的链，$X$为取代基或转移对象下标，$M^+$为金属反离子；$V_M,V_S$为单体/溶剂体积，$A,B$为两类互反应官能团。
\key{机理口诀：}给电子基（烷基、烷氧基）稳碳正离子$\to$阳离子；吸电子基（CN、COOR、双CN）稳碳负离子$\to$阴离子；苯基/二烯共轭可稳三类中心；乙烯基卤化物、St、MMA、AN、VAc、TFE常自由基；丙烯工业主要配位聚合。

\sect{1. 分子量、聚合度、分布}
结构单元是单体入链后的残余部分；重复单元是周期性最小单元。忽略端基：$M\approx XM_0$；Carothers 的 $\Xn$ 按结构单元计。$A_2+B_2$ 缩聚通常 $M_0$ 为重复单元质量一半：尼龙-66重复单元226，$M_0\approx113$；丁二醇--己二酸聚酯重复单元200，$M_0=100$。
\begin{gather*}
\Mn=\frac{\sum N_iM_i}{\sum N_i},\quad
\Mw=\frac{\sum N_iM_i^2}{\sum N_iM_i},\\
\PDI=\frac{\Mw}{\Mn}\ge1.
\end{gather*}
$\Mn$ 对低分子量敏感（端基、渗透压、GPC）；$\Mw$ 对高分子量敏感（光散射、GPC）。质量分数不能直接代入 $\Mn$ 的分子数公式。

\sect{2. 单体速判：作业全覆盖}
\subsect{适合机理}
\ce{CH2=CHCl}、\ce{CF2=CF2}：自由基；\ce{CH2=CHCN}：自由基/阴离子；\ce{CH2=C(CN)2}：阴离子；\ce{CH2=CHCH3}：配位为主，自由基低活性，阳离子条件可；\ce{CH2=C(CH3)2}、\ce{CH2=CHOCH3}：阳离子；\ce{CH2=CHPh}、\ce{CH2=CH-CH=CH2}：三类均可；\ce{CH2=C(CN)COOR}：阴离子/自由基。
\subsect{能否自由基聚合}
1,1-二苯基乙烯：难均聚（位阻大且自由基过稳）；\ce{ClHC=CHCl}、\ce{CF2=CFCl}：能；\ce{CH2=C(CH3)C2H5}：能但慢；\ce{CH3CH=CHCH3}、\ce{CH3CH=CHCOOCH3}：内烯烃难；MMA、醋酸乙烯酯能。判断顺序：电子效应$\to$位阻/取代度$\to$工业事实。

\sect{3. 链式 vs 逐步}
链式：有活性中心，引发/增长/终止分工，低转化率已有高分子量，分子量由 $\Ri,\kp,\kt$ 与链转移控制。逐步：官能团之间逐步反应，单体早期即消耗，只有 $p\to1$ 才有高分子量。实验判别：定时取样，用NMR/IR/滴定测转化率，用GPC/SEC/黏度/端基测 $\Mn$；低转化率已有高分子为链式，符合 $\Xn=1/(1-p)$ 且高分子量后期突增为逐步。

\sect{4. 自由基聚合历程}
$ I\xrightarrow{\kd}2R\cdot$；$R\cdot+M\xrightarrow{k_i}RM\cdot$；$P_n\cdot+M\xrightarrow{\kp}P_{n+1}\cdot$。偶合：$P_m\cdot+P_n\cdot\to P_{m+n}$，一条死链含两个引发剂残基；歧化：转移H生成两条死链，一条饱和一条不饱和，各含一个残基。AIBN先放 $N_2$ 得 $(CH_3)_2C(CN)\cdot$；BPO得 $\ce{PhCOO.}$，可脱羧为 $\ce{Ph.}$。阻聚：有诱导期；缓聚：速率下降但无明显诱导期；$t_{ind}\approx n[Z]/\Ri$。

\sect{5. 自由基动力学与自动加速}
理想假定：等活性、稳态、单体主要耗于增长、双基终止。
\begin{gather*}
\Ri=2f\kd[I],\quad
\Rp=-\frac{\dd[M]}{\dd t}=\kp[M][M\cdot],\\
\Rt=2\kt[M\cdot]^2.
\end{gather*}
稳态 $\Ri=\Rt$：
\[
[M\cdot]=\left(\frac{f\kd[I]}{\kt}\right)^{1/2},\quad
\Rp=\kp[M]\left(\frac{f\kd[I]}{\kt}\right)^{1/2}.
\]
故对 $[M]$ 一级、对 $[I]$ 半级；升温通常增速但降分子量。自动加速/凝胶效应：转化率高、黏度大、$k_t$ 受扩散限制下降，$[M\cdot]$ 和 $\Rp$ 突升；加溶剂/升温可推迟但可能增加链转移。

\sect{6. 动力学链长、终止比例、Mayo}
\[
\nu=\frac{\Rp}{\Ri},\qquad
\Xn=2\nu\ \text{(偶合)},\quad
\Xn=\nu\ \text{(歧化)}.
\]
若大分子平均含1.30个引发剂残基：先得偶合\emph{产物}分数 $x_p=0.30$、歧化产物0.70；但终止比例按反应物/活性链占比算，$f_c=0.3\times2/1.3=46.2\%$，$f_d=0.7/1.3=53.8\%$。链转移：$P_n\cdot+XA\xrightarrow{k_{tr}}P_nX+A\cdot$；$C_X=k_{tr,X}/k_p$。
\[
\frac1{\Xn}=\frac{\Ri}{2\Rp}+C_M+C_I\frac{[I]}{[M]}+C_S\frac{[S]}{[M]}+C_P\frac{[P]}{[M]}+\cdots.
\]
$C_X$ 越大，少量 $X$ 越能降分子量；转移后若新自由基可再引发，速率不一定显著降。
转移类型：单体(VC/VAc常显著限分子量)、引发剂、溶剂/硫醇调$\Mn$、大分子转移$\to$支化与PDI增大；新自由基弱则缓聚/阻聚。低转化分布：$\alpha=\Rp/(\Rp+\sum R_t+\sum R_{tr})$，歧化/转移 $\PDI=1+\alpha\le2$，偶合 $\PDI=(2+\alpha)/2\le1.5$；高转化、凝胶效应、支化会变宽。

\sect{7. 自由基数值模板}
二叔丁基过氧化物/St/苯：$\Ri=4.0\times10^{-11}$，$\Rp=1.5\times10^{-7}$，$[I]=0.01$。
\[
f\kd=\frac{\Ri}{2[I]}=2.0\times10^{-9}s^{-1},\quad \nu=\frac{\Rp}{\Ri}=3.75\times10^3.
\]
偶合无转移 $\Xn=7500$。苯溶剂浓度：1 L溶液含St 1 mol，$V_M=104.15/0.887=117.4$ mL，$V_S=882.6$ mL，$[S]=0.839\times882.6/78.11=9.49M$；代 Mayo 得 $\Xn\approx4.61\times10^3$。正丁硫醇降到 $\Mn=83200$：$\Xn=83200/104=800$，$1/800=1/7500+21[S]$，$[S]=5.32\times10^{-5}M$。

\sect{7a. 自由基共聚}
二元共聚：$r_1=k_{11}/k_{12}$，$r_2=k_{22}/k_{21}$；$f_i$ 是投料单体摩尔分数，$F_i$ 是瞬时共聚物组成：
\[
F_1=\frac{r_1f_1^2+f_1f_2}{r_1f_1^2+2f_1f_2+r_2f_2^2},\quad F_2=1-F_1.
\]
$r_1=r_2=1$ 理想无规，$F_1=f_1$；$r_1,r_2<1$ 交替倾向，均为0时严格交替；$r_1,r_2>1$ 倾向长同类序列/分别均聚；恒比点 $f_1^\ast=(1-r_2)/(2-r_1-r_2)$。马来酸酐难自聚但可与给电子单体交替共聚。

\sect{8. 引发剂/催化剂配对与条件}
BPO：自由基，St/VC/MMA；$t$-BuOOH/\ce{Fe^{2+}}：氧化还原自由基、低温；Na+萘/THF：电子转移阴离子，St、二氰基乙烯、MMA；\ce{H2SO4}：$\ce{H+}$ 阳离子，异丁烯、乙烯基醚、St；\ce{BF3+H2O}：$\ce{H+ [BF3OH]-}$，阳离子；$n$-BuLi：阴离子，St、二氰基乙烯、MMA。自由基溶剂惰性、温度按半衰期；阳离子需低温、无强亲核/碱性杂质；阴离子严格无水无氧无酸醇，溶剂常THF、苯、环己烷。

\sect{9. 阳离子聚合}
单体：异丁烯、乙烯基醚、能稳定正离子的St/二烯。引发剂：质子酸、Lewis酸+水/醇；反离子亲核性越弱越利于增长。增长中心为碳正离子/氧鎓离子对；溶剂极性增大可使离子对松散，活性升高。副反应：链转移、终止、重排；低温抑制副反应。
\begin{gather*}
\ce{BF3+H2O <=> H+ [BF3OH]-},\\
\ce{H+ + CH2=C(CH3)2 -> (CH3)3C+}.
\end{gather*}
速率式随机理变，考试重在影响方向：低温、弱亲核反离子、干燥、除杂提高可控性和分子量。

\sect{10. 阴离子/活性聚合}
单体：St、丁二烯、异戊二烯、MMA、AN、二氰基乙烯；引发剂：BuLi、萘钠、有机金属。理想活性聚合无终止/转移，活性链数恒定：
\begin{gather*}
\Rp=k_p[M][P^-],\quad
\ln\frac{[M]_0}{[M]}=k_p[P^-]t,\\
\Xn=\frac{[M]_0x}{[I]_{eff}}.
\end{gather*}
可顺序加料制嵌段，共轭二烯微结构受溶剂/反离子影响。萘钠/St：$1.0\times10^{-3}$ mol萘钠$\to5.0\times10^{-4}$ mol双阴离子链；2000 s转化50\%，$\Xn=2000$；达3000时 $x=0.75$，$t=2000\ln4/\ln2=4000$s。

\sect{11. RDRP与热力学}
传统自由基不可控因自由基浓度高、不可逆终止多。RDRP用“活性种$\rightleftharpoons$休眠种”降低自由基浓度并保留端基：NMP（氮氧自由基可逆捕获）、ATRP（卤代链/过渡金属可逆原子转移）、RAFT（二硫代酯/三硫代碳酸酯可逆加成--裂解）。聚合热力学：$\Delta G=\Delta H-T\Delta S$；乙烯基聚合多 $\Delta H<0,\Delta S<0$，$T_c=\Delta H/\Delta S$，低温热力学有利，高于 $T_c$ 可解聚；$T_c$ 用K，只判热力学不判速率。

\sect{12. 配位、立构、开环}
构型需断键改变，构象单键旋转即可。聚丙烯有全同/间同/无规；聚丁二烯有顺式、反式、1,2结构；立构规整性影响结晶度、密度、软化温度、强度、溶解性。Ziegler--Natta：\ce{TiCl3/TiCl4/VCl4}+\ce{AlR3/AlR2Cl}，单体先配位后插入金属--碳键，立体定向；链转移可向H$_2$/AlR$_3$/单体，H$_2$调分子量。ROP：环醚、内酯、内酰胺、环硅氧烷；三/四元环张力大易开环，五/六元环较稳定；尼龙-6来自己内酰胺。

\sect{13. 逐步聚合核心}
逐步聚合：官能团间反应，单体/低聚物/高聚物均可互反；每一步活性近似相同（等活性）。等当量线形：
\[
\Xn=\frac{N_0}{N}=\frac1{1-p},\quad \Xn=100\Rightarrow p=0.99;
\]
\[
\Xn=1000\Rightarrow p=0.999.
\]
不等当量或封端剂：
\[
\Xn=\frac{1+r}{1+r-2rp},\qquad (\Xn)_{max}=\frac{1+r}{1-r}\ (p=1).
\]
$r$ 由官能团数而非分子数直接决定；单官能杂质/单体损失会被高 $p$ 放大。

\newpage
\sect{14. 缩聚动力学、平衡、分布}
不可逆外加酸二级缩聚：$-\dd C/\dd t=kC^2$，$\Xn=1+C_0kt$。自催化按弱酸平衡 $[\ce{H+}]\simeq(K_aC)^{1/2}$，故 $-\dd C/\dd t=kC^{5/2}$，$\Xn^{3/2}=1+\frac32kC_0^{3/2}t$（不用讲义误写的三级）。可逆缩聚受平衡限制；提高分子量：高纯度、等当量、除水/醇/HCl、高真空/惰性气流/共沸、后期强化传质。最概然分布：
\begin{gather*}
N_x=N_0(1-p)^2p^{x-1},\\
\Xw=\frac{1+p}{1-p},\quad \PDI=1+p\to2.
\end{gather*}
$N_x$ 为聚合度 $x$ 的链数；逐步高转化分布天然接近2，不像活性阴离子接近1。

\sect{15. 逐步作业数值题}
尼龙-66加乙酸：重复单元226、结构单元 $M_0=113$，$\Mn=16000$，$\Xn=141.6$，$p=0.997$：$141.6=(1+r)/(1+r-2r\cdot0.997)$ 得 $r=0.9919$；胺基2 mol，总酸基 $2+x$，$r=2/(2+x)$，$x\approx0.016$ mol。丁二醇/己二酸：重复单元200、结构单元 $M_0=100$，$\Mn=5000$，$\Xn=50$，等当量 $p=0.980$。若丁二醇损失0.5\%，$r=0.995$，同 $p$ 下 $\Xn=44.53$、$\Mn\approx4.45\times10^3$；补加或预过量丁二醇。若羧基总2 mol中1\%为乙酸，$r=2.000/2.020=0.9901$，为 $\Xn=50$ 需 $p\approx0.985$。

\sect{16. 分子数法与封端剂}
对苯二甲酸/乙二醇/苯甲酸：羧基 $2.001$ mol，羟基 $2.000$ mol，$r=2.000/2.001=0.9995$；$p=0.95$，Carothers 得 $\Xn\approx19.9$。分子数法：初始分子 $N_0=2.001$ mol，反应羟基 $2.000\times0.95=1.900$ mol，每成1键分子数减1，$N=0.101$ mol，$\Xn=2.001/0.101\approx19.8$。封端剂题先画官能团账本：单官能酸只贡献一个羧基但不能继续延长链。

\sect{17. 非线形缩聚与凝胶}
支化：仍可溶熔；交联：三维网络，不溶不熔。凝胶点是首次出现无限网络的临界 $p$。Carothers：
\[
\bar f=\frac{\sum N_if_i}{\sum N_i},\qquad p_c=\frac2{\bar f}.
\]
若 $\bar f\le2$ 或 $p_c>1$，完全反应前通常不凝胶；Carothers 常偏大。Flory统计（等当量 $A_{f_A}+B$混合组）：
\[
f_{B,w}=\frac{\sum N_i f_i^2}{\sum N_i f_i},\quad p_c^2(f_A-1)(f_{B,w}-1)=1.
\]
邻苯二甲酸酐2.5 + 乙二醇1 + 丙三醇1：总官能团10，总分子4.5，$\bar f=2.222$，$p_c=0.900$；醇组 $f_{B,w}=13/5=2.6$，Flory $p_c=1/\sqrt{1.6}=0.791$。

\sect{18. 实施方法与典型聚合物}
自由基实施：本体（产品纯、放热难，油溶引发剂）、溶液（散热好但溶剂转移/后处理）、悬浮（单体珠滴小本体，油溶引发剂）、乳液（胶束/乳胶粒，水溶引发剂，速率快且可高分子量）。逐步实施：熔融缩聚（无溶剂，高真空脱小分子，PET/尼龙）、溶液缩聚（传热好）、界面缩聚（二酰氯+二胺，碱吸HCl）、固相缩聚（熔点以下提高PET/尼龙分子量）。典型：PET、PTT、尼龙-66、尼龙-6、聚碳酸酯、聚氨酯、环氧、酚醛、脲醛、不饱和聚酯。

\sect{19. 物理量实验来源}
$p$：端基滴定、酸值/羟值、IR峰、NMR积分；双键转化率可用NMR/IR/GC。$\Mn$：端基分析、膜渗透压、GPC校准；端基法适合分子量不太高且端基定量。$\Mw$：光散射、GPC，对高分子量尾部敏感。$\Rp\approx[M]_0\dd x/\dd t$（体积近似不变）。$k_p,k_t$ 若题给 $[M\cdot]$ 可由 $\Rp=\kp[M][M\cdot]$、$\Rt=2\kt[M\cdot]^2$ 反推。

\sect{20. 公式适用条件}
自由基理想式若有自动加速、强阻聚、强链转移，就不能机械套；$\Ri=2f\kd[I]$ 只适合一级分解且每分子给两个自由基，光/氧化还原体系可能另给 $\Ri$。$\Xn=1/(1-p)$ 只适合等当量线形逐步上限；有封端剂、单体挥发、环化、不等活性需修正。凝胶公式默认等反应性、无分子内环化；实际环化推迟凝胶。活性阴离子公式默认无终止、无转移、活性中心数恒定。

\sect{21. 短答标准句}
半级原因：稳态下自由基生成与双基终止平衡，$[M\cdot]\propto[I]^{1/2}$，故 $\Rp$ 对 $[I]$ 半级。高 $p$ 要求：逐步聚合每形成一个键只使分子数减1，$p$ 接近1时少量未反应官能团决定分子量。链转移：原链提前停止，新自由基可能继续增长，所以主要降分子量，速率未必同比例降。凝胶：从一个支化点平均产生超过一个新支化点时，网络无限扩展，溶胶转凝胶。RDRP可控：休眠/活性快速交换，低自由基浓度降低终止并保留端基。

\sect{22. 最后防错清单}
重复单元质量扣小分子；$p$ 是官能团反应程度，不是分子转化率；$r$ 先数官能团且取 $\le1$；封端剂显著限 $\Xn$；$T_c$ 用K；密度题先算体积和浓度；$\PDI<1$ 必错；$\Xn<1$ 多为定义/单位错；三官能/四官能物增多，$\bar f,f_w$ 增大，$p_c$ 应降低；$[I]$ 增大，$\Rp$ 升但 $\nu$ 和分子量降；内烯烃/多取代烯烃自由基均聚常难；离子聚合最怕相反性质杂质（阳离子怕亲核/碱，阴离子怕酸/水/醇/氧）。

\sect{23. 常见缩写}
St 苯乙烯；MMA 甲基丙烯酸甲酯；AN 丙烯腈；VC 氯乙烯；VAc 醋酸乙烯酯；TFE 四氟乙烯；BPO 过氧化苯甲酰；AIBN 偶氮二异丁腈；NMP 氮氧自由基聚合；ATRP 原子转移自由基聚合；RAFT 可逆加成--裂解链转移；ROP 开环聚合；SEC/GPC 凝胶渗透/尺寸排阻色谱；Z--N 齐格勒--纳塔。


\sect{24. 链式反应式书写模板}
自由基：$R\cdot+\ce{CH2=CHX}->R\ce{-CH2-\dot{C}HX}$；连续增长写 $P_n\cdot+\ce{CH2=CHX}->P_{n+1}\cdot$。阳离子：$\ce{H+ + CH2=CHX -> CH3-CHX+}$，必须带反离子如 $\ce{HSO4-}$、$\ce{BF3OH-}$；增长为碳正离子加双键。阴离子：$R^-M^++\ce{CH2=CHX}->R\ce{-CH2-CHX^-}M^+$，要标明金属反离子。配位聚合短答必须写“单体配位到金属中心、插入金属--碳键、立体定向”。

\sect{25. 单体电子/位阻细化}
自由基聚合既要求生成的增长自由基有适当稳定性，也要求双键足够活泼、位阻不太大。自由基过稳会使继续加成动力不足，如1,1-二苯基乙烯；内烯烃和三/四取代烯烃因位阻和热力学不利常难均聚。氟代烯烃是例外，\ce{CF2=CF2}、\ce{CF2=CFCl} 可自由基聚合。阳离子聚合要求正离子被烷基、烷氧基或苯基稳定；阴离子聚合要求负离子被氰基、羰基、酯基、苯基或共轭体系稳定。

\sect{26. 三类链式聚合比较}
自由基：单体适用面广，氧阻聚，水/杂质容忍度中等；引发剂为偶氮、过氧化物、氧化还原；终止为偶合/歧化，链转移常见。阳离子：给电子单体，质子酸或Lewis酸，怕水/醇/胺等亲核或碱性杂质，常低温，转移多。阴离子：吸电子/共轭单体，BuLi/萘钠，怕水氧酸醇，理想时活性聚合、窄分布、可嵌段。

\sect{27. 聚合热力学补充}
聚合可行需同时满足热力学和动力学。$\Delta H$ 越负越有利，但共轭取代基稳定单体会降低聚合放热；$\Delta S<0$ 来自多个单体变一条链后平动/转动自由度下降。接近 $T_c$ 时聚合与解聚平衡显著；工业上还需考虑散热、黏度和副反应。开环聚合还受环张力和构象熵影响，小环通常易开环，大环可由熵因素驱动。

\sect{28. RDRP三法辨析}
NMP中休眠链末端与氮氧自由基形成可逆烷氧胺，适合St等；ATRP中 $P-X$ 与低价金属络合物可逆卤原子转移，保留卤素端基，适合嵌段/接枝；RAFT中硫代羰基硫化合物控制可逆加成--裂解，单体适用广。共同结果：$\Mn$ 随转化率近线性增长，PDI降低，端基保留；不是“完全无终止”，而是终止概率显著降低。

\sect{29. 逐步聚合单体官能度}
单官能物不能成链，只能封端/调分子量；$A_2+B_2$ 或 $AB$ 生成线形聚合物；$A_2+B_f$、$A_f+B_2$ 且 $f>2$ 可支化/交联；$AB_f$ 常生成超支化结构。官能度计数按“可反应基团数”而不是分子名称：二酸有2个羧基，二醇有2个羟基，丙三醇有3个羟基，酸酐在缩聚中常按二官能酸处理。

\sect{30. 提高逐步聚合分子量措施}
保持单体高纯、严格等当量；移除水、醇、HCl等小分子；使用高真空、惰性气流、共沸带水；选择适当催化剂；后期提高温度但避免降解；提高搅拌和传质；避免单官能杂质、单体挥发和环化。若已发生二醇挥发，应补加二醇或预先略过量，而不是只延长时间。

\sect{31. 逐步计算流程}
1 算重复单元 $M_0$，目标 $\Xn=\Mn/M_0$；2 判断等当量：等当量用 $\Xn=1/(1-p)$；3 不等当量/封端剂先列官能团总数，取 $r\le1$，用 $(1+r)/(1+r-2rp)$；4 也可分子数守恒校验；5 若涉及凝胶，先判断是否有 $f>2$ 单体，再选 Carothers 或 Flory。

\sect{32. 凝胶题流程}
Carothers适合快速估算：算所有有效分子的 $\bar f$，$p_c=2/\bar f$。若官能团不等量，应先修正为可进入网络的有效组成。Flory适合 $A_2+B_2+B_3$ 等统计题：先确认A/B等当量，算支化侧的 $f_w=\sum N_if_i^2/\sum N_if_i$，再代 $p_c^2(f_A-1)(f_w-1)=1$。高官能度单体因平方加权影响很大。

\sect{33. 分布与性能联系}
分子量升高通常提高强度、韧性、熔体黏度和加工难度；低分子尾部会降低力学性能，高分子尾部会显著提高 $\Mw$ 和熔体弹性。活性阴离子/RDRP分布窄，适合结构明确嵌段；普通自由基和逐步缩聚分布较宽。逐步理想高转化 PDI 接近2，不代表有副反应。

\sect{34. 实验判别补充}
链式聚合的单体转化率随时间逐渐上升，早期聚合物分子量已接近最终量级；逐步聚合早期多为二聚/三聚/低聚物，$\Mn$ 后期随 $p$ 急剧上升。可结合黏度变化、GPC曲线、端基浓度、官能团峰衰减判断。未知单体实验应固定温度/引发剂，定时淬灭取样，避免后反应造成假象。

\sect{35. 典型聚合物结构来源}
PVC来自氯乙烯自由基；PAN来自丙烯腈；PMMA来自MMA；PS来自St；PTFE来自TFE；聚异丁烯和丁基橡胶常由阳离子聚合；SBS可由活性阴离子顺序加料；聚乙烯/聚丙烯可由配位聚合；PET由二酸/二酯与乙二醇缩聚；尼龙-66由己二酸+己二胺；尼龙-6由己内酰胺开环；PC由双酚A与光气或碳酸二苯酯。

\sect{36. 方程单位快查}
$k_d$ 为 $s^{-1}$；$k_p,k_t,k_{tr}$ 多为 $L\,mol^{-1}s^{-1}$；$\Ri,\Rp,\Rt$ 为浓度/时间；$C_X,f,p,r,\nu,\Xn,\PDI$ 无量纲；$\Delta H$ 常用 $J\,mol^{-1}$，$\Delta S$ 用 $J\,mol^{-1}K^{-1}$，所以 $T_c$ 为K。若结果 $p>1$、PDI$<1$、$\Xn<1$，优先查单位和定义。

\sect{37. 作业中最容易漏写}
AIBN题要写 $N_2$；BPO题要写酰氧自由基及可脱羧；离子聚合式要写反离子；萘钠题要注意2 mol萘钠对应1 mol双阴离子链；苯溶剂题要由密度算 $[S]$；尼龙-66题乙酸是单官能酸、加入到酸基总数；对苯二甲酸/苯甲酸题两种算法结果略有差别但应接近；凝胶题Carothers与Flory结果不同且Flory更低。

\end{multicols}
\end{dense}
\end{document}
